\homework{2}{Mon 29 Aug 2022 14:24}{Week 2 Due Sep 9}

I use the following lemma to prove Exercise 3.J.
\begin{lemma}[pigeonhole principle]
	Let $Y$ be a finite set and $X$ be finite with cardinality greater than $|Y|$. Then for any function $f: X \to  Y$, there exists $x_1, x_2 \in X$ such that $f(x_1) = f(x_2)$.
\end{lemma}
\begin{proof}
Denote $A_y$ by the set $\{x \in X: f(x) = y\} $. We must have $|A_y| < 2$ for all $y \in Y$, otherwise there will be more than one element from $X$ mapped to the same element in $Y$. Hence $0 \le |A_y| \le  1$. We have
\[
\left| \bigcup_{y \in Y} A_y \right| = \sum_{y \in Y } |A_y| \le \sum_{y \in Y} 1 =  |Y|
.\] 
The first equality comes from the fact that $A_y$ are all disjoint and finite (since elements mapped to different targets must be different).

But for any $x \in X$, it is mapped to some element $y = f(x)$. In other words, $x \in A_{f(x)} \subset \bigcup_{y \in Y} A_y$. In the other direction, for any $x \in \bigcup_{y \in Y} A_y$, it is some element in $X$ by definition. We have shown that $X = \bigcup_{y \in Y} A_y$ and therefore $|X| = \left| \bigcup_{y \in  Y} A_y \right| \le |Y|$ which is a contradiction to the assumption that $|X| > |Y|$. Therefore there must exist some $y$ such that  $|A_y| \ge 2$, which is exactly what we want to prove.
\end{proof}

I prove this exercise in extra, because I used it in Exercise 7.I.

7.I. Show that the collection of "right hand" end points in $\boldsymbol{F}$ is denumerable. Show that if all these end points are deleted from $\boldsymbol{F}$, then what remains can be put onto one-one correspondence with all of $[0,1)$. Conclude that the set $\boldsymbol{F}$ is not countable.
\begin{proof}
Each "right hand" end point is in the form $\frac{k+1}{3^n}$ where $k, n \in \mathbb{N}$. Therefore it is a rational number, and the set of rational numbers $\mathbb{Q}$ is denumerable. Because the set of right hand end points is a subset of rational numbers, it is denumerable.

Now we construct cascaded segmentation of the half-open interval $[0,1)$. For each nonnegative integer $n$ and $0\le k\le n$, define $A_{nk}$ to be \[
A_{nk} = [\frac{k}{2^n}, \frac{k+1}{2^n}) 
.\] 

We use induction to verify the following properties:
\begin{enumerate}
	\item for all $n$ there are exactly $2^n$ disjoint sets of the form $A_{nk }$.
	\item (disjoint segmentation) for all $n$ and all distinct  $k, l$, 
		\[
		A_{nk}\cap A_{nl} = \emptyset 
		.\] 
	\item (filling) for all $n$, \[
\bigcup_{0\le k\le n} A_{nk} = [0,1)
.\] 

\item (branch split) for any $n, k$, $A_{n+1, l} \subset A_{nk}$ if and only if \[
2k \le l \le 2k+1
.\] 
\item (cascading property) whenver $x \in A_{nk}$ and $x \in A_{ml}$ where $n \le  m$, we have $A_{ml} \subset  A_{nk}$.
\end{enumerate}

Now for each $n$, we can construct a bijection $G:\bigcup_{0\le k\le n} A_{nk} \to F_n^*$, where $F_n$ is as defined in the textbook, and $F_n^*$ is  $F_n$ with right endpoints removed. 
There are exactly $2^n$ disjoint half-open intervals $A_{nk }$. There are also exactly $2^n$ disjoint half-open intervals in $F_n^*$. We map the $k$-th interval $A_{nk }$ to the $k$-th interval in $F_n^*$. Note that any two half-open intervals on $\mathbb{R}$ have a bijection between them.

Now we show that there exists a bijection $\mathcal{F}:[0,1) \to F^*$ where $F^*$ is the Cantor set with right endpoints removed. Take any $x \in [0,1)$. For any $n \in \mathbb{N}$, let $A_{nk}$ be the unique half-open interval in the construction that contains $x$. Now $G$ maps $A_{nk }$ to some half-interval $B_n$ in the construction of $F_n^*$. 
Now $B_{n+1}\subset B_n$ and  $B_n \neq \emptyset $. Denote $b_n$ as the left endpoint of $B_n$. Now $y = \operatorname{sup}_{n \in \mathbb{N}} b_n $ is the unique element in $\bigcap_{n \in \mathbb{N}} B_n \subset F^*$.


Such map is injective: if $x_1<x_2$ are distinct point in $[0,1)$, then there exists $n$ such that $\frac{1}{2^n}\le \frac{1}{n} < x_2-x_1$. Then $x_1 \in A_{nk }$ and $x_2 \in A_{nl }$ where $n\neq l$, hence they are mapped to disjoint half-open intervals $B_1$ and $B_2$ in $F_n^*$. Now $\mathcal{F}(x_1) \in B_1$ and $\mathcal{F}(x_2) = B_2$, so they must not be equal. 

Such map is surjective: if $y \in F^*$, then for any $n \in \mathbb{N}$, $y \in B_n \subset F_n^*$ where $B_n$ is a half-open interval in the construction of $F_n^*$. $\mathcal{F}^{-1}$ then maps $B_n$ to some $A_n$. Since $B_{n+1} \subset B_n$, $A_{n+1} \subset A_n$. Let $a_n$ be the endpoint of $A_n$, and $b_n$ the left endpoint of $B_n$. Now $x = \operatorname{sup}_{n \in \mathbb{N}}a_n$ is the unique element in $\bigcap_{n \in \mathbb{N} }A_n $. By definition $\mathcal{F}(x) = \operatorname{sup} _{n \in \mathbb{N}} b_n$. This is the unique element in $\bigcap_{n \in \mathbb{N}} B_n$. Since $y \in \bigcap_{n \in \mathbb{N}} B_n$, $\mathcal{F}(x) = y$.
\end{proof}
          

6.M, 7.F, 7.G, 7.J, 7.L, 3.F, 3.J

\begin{enumerate}
	\item 6.M. Modify the argument given in Theorem $6.8$ to show that if $a>0$, then the number
\[
b=\sup \left\{y \in \boldsymbol{R}: 0 \leq y, \quad y^{2} \leq a\right\}
\]
exists and has the property that $b^{2}=a$. This number will be denoted by $\sqrt{a}$ or $a^{1 / 2}$, and is called the positive square root of $a$.
\begin{proof}
	Denote the set by $S$. We first show that the supremum exists. Clearly $0 \in S$, so the set is nonempty. Net we show that the set is bounded from above. We split into cases where $x\le 1$ and $x>1$. When  $x\le 1$, we claim that the set is bounded above by $1$. Suppose not, then there exists some $s \in S$ such that $s >1$, hence $s^2 >1$ but we have $s^2 \le a \le 1$, a contradiction. When $x>1$, we claim that the set is bounded above by $a$. Suppose not, then there exists some $s \in S$ such that $s > a$, hence $s^2 > a^2$ but we have $s^2 \le  a < a^2$ (the second inequality is due to $a>1$), a contradiction. In both cases the set is bounded above. By supremum property of $\mathbb{R}$, $b = \operatorname{sup} S$ exists.

	Now we show $b^2 = a$. Suppose not, then by trichotomy property of real ordering, we have either $b^2 < a$ or $b^2 > a$. When $b^2 < a$, we take $n \in \mathbb{N}$, such that $\frac{1}{n}<\frac{a-b^2}{2b+1}$ (using Corollary 6.7 (b) of Archimedean Property). In this case \[
		(b+\frac{1}{n})^2 = b^2 + \frac{1}{n}\left( 2b + \frac{1}{n} \right) \le b^2 + \frac{1}{n}\left( 2b + 1 \right) < b^2 + (a-b^2) = a
	.\] 
This means $b + \frac{1}{n} \in S$, contrary to the fact that $b$ is an upper bound of $S$.

When $b^2>a$, we take $n \in \mathbb{N}$ such that $\frac{1}{n} < \frac{b^2 - a}{2b}$ (such $n$ exists by Corollary 6.7(b) of the Archimedean Property). Since $b = \operatorname{sup} S$, there exists an $s_0\in S$ such that $b - \frac{1}{n} <s_0$. But this implies that
\[
a < b^2 - 2 \frac{b}{n} < b^2 - 2 \frac{b}{n} + \frac{1}{n^2} = \left( n-\frac{1}{n} \right) ^2 < s_0^2
.\] 
But this means $s_0^2 > a$, contrary to the fact that $s_0 \in S$.

We have excluded the possibilities of $b^2<a$ and $b^2>a$, so the only possibility left is $b^2 = a$.
\end{proof}

	\item 7.F. Let $J_{n}=(0,1 / n)$ for $n \in \mathbf{N}$. Show that this sequence of intervals is nested, but that there is no common point.
\begin{proof}
	We first show the sequence is nested. Fix some $n \in \mathbb{N}$. It suffices to show that $J_{n+1} \subset J_n$. Take some $x \in J_{n+1}$, then by definition we have $0<x<\frac{1}{n+1}$. But since $\frac{1}{n+1}<\frac{1}{n}$, we have $0<x<\frac{1}{n}$, hence by definition again we have $x \in J_n$.

	Now we show there is no common point. Fix any $x \in J_1 = (0,1)$. By Corollary 6.7(b), there exists some $n \in \mathbb{N}$ such that $0<\frac{1}{n}<x$. In other words, $x \not\in J_n$. Therefore, $\bigcap_{n \in \mathbb{N}} J_n = \emptyset $ since each element is not in at least one of the sets.
\end{proof}
	\item 7.G. If $I_{n}=\left[a_{n}, b_{n}\right], n \in \mathbf{N}$, is a nested sequence of closed cells, show that
\[
a_{1} \leq a_{2} \leq \cdots \leq a_{n} \leq \cdots \leq b_{m} \leq \cdots \leq b_{2} \leq b_{1} .
\]
If we put $\xi=\sup \left\{a_{n}: n \in \mathbf{N}\right\}$ and $\eta=\inf \left\{b_{m}: m \in \mathbf{N}\right\}$, show that $[\xi, \eta]=\bigcap_{n \in N} I_{n}$.
\begin{proof}
	Take $n, m \in \mathbb{N}$ where $n\le m$. By assumption $I_m \subset  I_n$. We need to show that $a_n \le a_m$ and $b_m\le b_n$. By definition of closed intervals, for all  $x \in \mathbb{R}$, $x \in I_n \iff a_n \le  x \le  b_n$ and $x \in I_m \iff a_m \le  x \le  b_m$. By set inclusion, $x \in I_m \implies x \in I_n$. Therefore \[
	a_m \le  x \le  b_m \implies a_n \le  x \le  b_n
	.\] 
Suppose that $a_m < a_n$. Then there exists some $x$ such that $a_m < x < x_n$, such that $a_n \le x\le b_n$ is not true, therefore the implication is not true for the specific $x$, a contradiction to the universally quantified statement. By trichotomy of real order, we have $a_m \ge a_n$. By the same reasoning we have $b_m \le  b_n$.

\begin{correction}
	We still need to verify that $a_n\le  b_m$ for all $n, m \in  \mathbb{N}$. To do this, take $N = \max(n,m)$. Now \[
		a_n \le a_N \le b_N \le b_m
	.\] 
\end{correction}
Now we show that $[\xi, \eta]=\bigcap_{n \in N} I_{n}$. We first show that $[\xi, \eta]\subset \bigcap_{n \in N} I_{n}$. Take $x \in  \left[ \xi, \eta \right] $, and take an arbitrary $n \in \mathbb{N}$. It suffices to show that $x \in I_n$. Since $\xi$ is a supremum, we must have $\xi \ge a_n$. Therefore $x \ge  \xi \ge  a_n$. Similarly we have $x \le  b_n$. By definition of $I_n$, we have $x \in  I_n$.

Next we show that $[\xi, \eta]\supset \bigcap_{n \in N} I_{n}$. Take $x \in \bigcap_{n \in N} I_{n}$. Then $x \in I_n$ for all $n \in \mathbb{N}$. Suppose that $x < \xi$. Then there exists some $m \in \mathbb{N}$ such that $x<a_m \le  \xi$. Clearly $x \in I_m$ is not true, a contradiction. Therefore we must have $x \ge \xi$. By the same reasoning we have $x\le \eta$, so $x \in  \left[ \xi, \eta \right] $.
\end{proof}
	\item 7.J. Every open interval $(a, b)$ which contains a point of $\boldsymbol{F}$ also contains an entire "middle third"' set which belongs to $\mathscr{C}(\boldsymbol{F})$. Hence $\boldsymbol{F}$ does not contain any non-void open interval.
\begin{proof}
	Fix some interval $(a,b) \subset \mathbb{R}$ and some point $x \in \mathbf{F}\cap (a,b)$. Since $a<x<b$, we have $b-x>0$ and $x-a>0$. We can choose some $n \in \mathbb{N}$ such that\[
	3^n > \max(\frac{1}{x-a}, \frac{1}{b-x})
	.\] 
	To justify this, note that by Bernoulli's Inequality\[
	3^n = (1+2)^n \ge  1 + 2n > n
	\]
	and by Archimedean property we can choose $n$ such that  $n> \max(\frac{1}{x-a}, \frac{1}{b-x}) $.

	By Corollary 6.7(c) of the Archimedean property, there exists $k \in \mathbb{N}$ such that $k \le  3^n x < k+1$. Since $3^n x - 3^n a > 1$, we have $k > 3^n x -1 > 3^n a$. Since $3^nb - 3^nx > 1$, we have $k+1 \le 3^n x + 1<3^n b$. Therefore $\left[ \frac{k}{3^n}, \frac{k+1}{3^n} \right]  \subset  \left( a,b \right) $.

	By definition and construction, $x \in \left[ \frac{k}{3^n}, \frac{k+1}{3^n} \right] \subset F_n$. $F_{n+1}$ takes the middle third from this interval, which is $\left[ \frac{3k}{3^{n+1}}, \frac{3k+1}{3^{n+1}} \right]\cup \left[ \frac{3k+2}{3^{n+1}}, \frac{3k+3}{3^{n+1}} \right] \subset F_{n+1}$. The middle-third $M = \left(  \frac{3k+1}{3^{n+1}}, \frac{3k+2}{3^{n+1}}  \right)$ satisfies the properties $M \subset F_n $ and $M \subset (a,b)$.
\end{proof}

	\item 7.L. Show that $\boldsymbol{F}$ is not the union of a countable collection of closed intervals.
\begin{proof}
	Suppose $\mathbf{F} = \bigcup_{n \in \mathbb{N}} I_n$ where each $I_n = [a_n, b_n]$. Suppose there exists some $m \in \mathbb{N}$ such that $a_m < b_m$. Then the open interval $(a_m, b_m)$ is nonempty, and has an nonempty intersection with $\mathbf{F}$. By Exercise 7.J, the interval $(a_m, b_m)$ contains, a "middle-third" set. This contradicts the fact that $[a_n, b_n] \subset  \mathbf{F}$. 

	Therefore $a_n\ge b_n$ for all $n \in \mathbb{N}$. If $a_n>b _n$ then the interval is empty, so we have $I_n = \{x_n\} \subset \mathbf{F}$ for at most countable number of $I_n$, and the rest of them are empty. This is to say that the set $\mathbf{F}$ is at most countable. But we have concluded in Exercise 7.I that  $\mathbf{F}$ is not countable.
\end{proof}

	\item 3.F. Use the fact that every infinite set has a denumerable subset to show that every infinite set can be put into one-one correspondence with a proper subset of itself.
\begin{proof}
Consider an infinite set $X$ and a denumerable subset  $A\subset X$. Denumerate $A$ by $a_1, a_2, \ldots, a_i, \ldots$, and define the map \begin{align*}
	F: X &\longrightarrow X\setminus \{a_1\}  \\
	x &\longmapsto F(x) = \begin{cases}
		x & x \not\in A\\
		a_{i+1} & x = a_i \in A
	\end{cases}
.\end{align*}
$F$ is a function because for every $x$ either $x \not\in A$, or $x \in A$, in this case there must exist $i \in \mathbb{N}$ such that $x = a_i$. $F$ is injective since for any $x_1 \neq  x_2$, either they both  $\not\in  A$, in this case $F(x_1) = x_1 \neq  x_2 = F(x_2)$; or they are both $\in A$, in this case $F(x_1) = F(a_i) = a_{i+1} \neq a_{j+1} = F(a_j) = F(x_2)$ ; or one of them is in $A$ and the other is not, in this case the claim is obvious.

$F$ is surjective because for any $x \in X\setminus A$, its preimage is itself; and for any $x = a_i \in A\setminus \{a_1\} $ where $i \neq 1$, its preimage is $a_{i -1} \in A\setminus \{a_1\} $.
\end{proof}


\item 3.J. Prove that $\mathbf{N}$ cannot be put into one-one correspondence with any initial segment of $\mathbf{N}$.
\begin{proof}
Suppose for the sake of contradiction that there exists some $n \in \mathbb{N}$ such that there is a bijection $f: \mathbb{N} \to  S_n$. Hence in particular $f$ is injective. Then by restricting the domain, $f|_{S_{n+1}} : S_{n+1} \to  S_n$ is injective. However this cannot be true since by the pigeonhole principle there must exist two elements in $S_{n+1}$ that are mapped to the some same element in $S_n$.
\end{proof}

\end{enumerate}
